\documentclass{article}

\usepackage{braket}
\usepackage{amsmath}

\begin{document}
    \begin{align*}
        \ket{\psi_0} &= \ket{x} \\ \\ \\
        \ket{\psi_1} &= H_1\ket{\psi_0} \\
        \ket{\psi_1} &= \frac{1}{\sqrt{2}}\ket{0} + (-1)^x \frac{1}{\sqrt{2}}\ket{1} \\
        \ket{\psi_1} &= \frac{1}{\sqrt{2}}(\ket{0} + (-1)^x \ket{1}) \\ \\ \\
        \ket{\psi_2} &= X_1\frac{1}{\sqrt{2}}(\ket{0} + (-1)^x\ket{1}) \\
        \ket{\psi_2} &= \frac{1}{\sqrt{2}}(X_2\ket{0} + (-1)^xX_2\ket{1}) \\
        \ket{\psi_2} &= \frac{1}{\sqrt{2}}\ket{1} + (-1)^x\frac{1}{\sqrt{2}}\ket{0} \\ \\ \\
        \ket{\psi_3} &= H_1\frac{1}{\sqrt{2}}(\ket{1} + (-1)^x\ket{0}) \\
        \ket{\psi_3} &= \frac{1}{\sqrt{2}}(H_3\ket{1} + (-1)^xH_3\ket{0}) \\
        \ket{\psi_3} &= \frac{1}{2}(\ket{0} - \ket{1} + (-1)^x (\ket{0} + \ket{1})) \\
    \end{align*}

Essayons maintenant d'evaluer le résultat si $\ket{x} = \ket{0}$ :\\
    \begin{align*}
        \ket{\psi_3} &= \frac{1}{2}(\ket{0} - \ket{1} + (-1)^0 \ket{0} + \ket{1}) \\
        \ket{\psi_3} &= \frac{1}{2}(\ket{0} - \ket{1} + \ket{0} + \ket{1}) \\
        \ket{\psi_3} &= \frac{1}{2}(2\cdot\ket{0}) \\
        \ket{\psi_3} &= \ket{0} \\
    \end{align*}
Ensuite, si $\ket{x} = \ket{0}$ : \\
    \begin{align*}
        \ket{\psi_3} &= \frac{1}{2}(\ket{0} - \ket{1} + (-1)^1 \ket{0} + \ket{1}) \\
        \ket{\psi_3} &= \frac{1}{2}(\ket{0} - \ket{1} - \ket{0} - \ket{1}) \\
        \ket{\psi_3} &= \frac{1}{2}(-2\cdot\ket{0}) \\
        \ket{\psi_3} &= -\ket{1} \\
    \end{align*}
Donc, la porte $HXH$ appliqué à $\ket{x}$ donne :
\[
HXH = 
    \left\{
        \begin{aligned}
            &\ket{0} \mapsto \ket{0}\\
            &\ket{1} \mapsto -\ket{1}
        \end{aligned}
    \right .
\] \\

Ce qui a le même effet que la porte $Z$.

\end{document}